%% TRANSCRIBED (PDF 50-55)
%% Chapter II opening + Sec. 1 Preliminaries (printed II-1 .. II-6).
%% The final Exercise of Sec. 1.3 ends at the top of II-6 (PDF 55),
%% immediately before Sec. 2 (transcribed in ch2_s2.tex).

% ---- chapter opening (printed II-1) ----

Throughout this chapter, $K$ denotes an \emph{algebraic number field}.
We associate to $K$ a projective family $(S_m)$ of commutative algebraic
groups over $\Q$, and we show that each $S_m$ gives rise to a strictly
compatible system of rational \textit{l}-adic representations of $K$.

In the next chapter, we shall see that all ``locally algebraic'' abelian
rational representations are of the form described here.

\sechead{1}{Preliminaries}{PRELIMINARIES}

\subsechead{1.1}{The torus $T$}

Let $T = \mathrm{R}_{K/\Q}(\mathbf{G}_{m/K})$ be the algebraic group over
$\Q$, obtained from the multiplicative group $\Gm$ by restriction of
scalars from $K$ to $\Q$, cf.~Weil~\cite{weil-adeles}, \S1.3.  If $A$ is a
commutative $\Q$-algebra, the points of $T$ with values in $A$ form by
definition the multiplicative group $(K \tens_\Q A)^*$ of invertible
elements of $K \tens_\Q A$.  In particular, $T(\Q) = K^*$.  If
$d = [K : \Q]$, the group $T$ is a torus of dimension $d$; this means that
the group $T_{/\bar{\Q}} = T \times_\Q \bar{\Q}$ obtained from $T$ by
extending the scalars from $\Q$ to $\bar{\Q}$, is isomorphic to
$\mathbf{G}_{m/\bar{\Q}} \times \ldots \times \mathbf{G}_{m/\bar{\Q}}$
($d$ times).  More precisely, let $\Gamma$ be the set of embeddings of $K$
into $\bar{\Q}$; each $\sigma \in \Gamma$ extends to a homomorphism
$K \tens_\Q \bar{\Q} \to \bar{\Q}$, hence defines a morphism
$[\sigma]\colon T_{/\bar{\Q}} \to \mathbf{G}_{m/\bar{\Q}}$.  The
collection of all $[\sigma]$'s gives the isomorphism
$T_{/\bar{\Q}} \to \mathbf{G}_{m/\bar{\Q}} \times \ldots \times
\mathbf{G}_{m/\bar{\Q}}$.  Moreover, the $[\sigma]$'s form a basis of the
\emph{character group} $X(T) = \Hom_{\bar{\Q}}(T_{/\bar{\Q}},
\mathbf{G}_{m/\bar{\Q}})$ of $T$.  Note that the Galois group
$\Gal(\bar{\Q}/\Q)$ acts in a natural way on $X(T)$, viz.~by permuting the
$[\sigma]$'s.  (For the dictionary between tori and Galois modules, see
for instance T.~Ono~\cite{ono-tori}.)

\subsechead{1.2}{Cutting down $T$}

Let $E$ be a subgroup of $K^* = T(\Q)$ and let $\bar{E}$ be the
\emph{Zariski closure} of $E$ in $T$.  Using the formula
$\bar{E} \times \bar{E} = \overline{E \times E}$, one sees that $\bar{E}$
is an \emph{algebraic subgroup} of $T$.  Let $T_E$ be the quotient group
$T/\bar{E}$; then $T_E$ is also a torus over $\Q$.  Its character group
$X_E = X(T_E)$ is the subgroup of $X = X(T)$ consisting of those
characters which take the value $1$ on $E$.  If
$\lambda = \prod\limits_{\sigma \in \Gamma} [\sigma]^{n_\sigma}$ denotes a
character of $T$, then $X_E$ is the subgroup of those $\lambda \in X$ for
which $\prod \sigma(x)^{n_\sigma} = 1$, for all $x \in E$.

\runhead{Exercise}
\begin{alist}
\item[\emph{a.}] Let $K$ be quadratic over $\Q$, so that $\dim T = 2$.
Let $E$ be the group of units of $K$.  Show that $T_E$ is of dimension
$2$ (resp.~$1$) if $K$ is imaginary (resp.~real).
\item[\emph{b.}] Take for $K$ a cubic field with one real place and one
complex one, and let again $E$ be its group of units (of rank~1).  Show
that $\dim T = 3$ and $\dim T_E = 1$.
\end{alist}

(For more examples, see \S\ref{sec:II.3.3}.)

\subsechead{1.3}{Enlarging groups}

Let $k$ be a field and $A$ a commutative algebraic group over $k$.  Let
\begin{equation}
\tag{$*$}\label{eq:II.1.3.star}
0 \longrightarrow Y_1 \longrightarrow Y_2 \longrightarrow Y_3
\longrightarrow 0
\end{equation}
an exact sequence of (abstract) commutative groups, with $Y_3$
\emph{finite}.  Let
\[
\epsilon : Y_1 \longrightarrow A(k)
\]
be a homomorphism of $Y_1$ into the group of $k$-rational points of $A$.
We intend to construct an algebraic group $B$, together with a morphism of
algebraic groups $A \to B$ and a homomorphism of $Y_2$ into $B(k)$ such
that,
\begin{alist}
\item \emph{the diagram}
\[
\begin{array}{ccc}
Y_1 & \longrightarrow & A(k) \\
\downarrow & & \downarrow \\
Y_2 & \longrightarrow & B(k)
\end{array}
\]
\emph{is commutative},
\item $B$ \emph{is ``universal'' with respect to} (a).
\end{alist}

The universality of $B$ means that, for any algebraic group $B'$ over $k$
and morphisms $A \to B'$, $Y_2 \to B'(k)$ such that (a) is true (with $B$
replaced by $B'$), there exists a unique algebraic morphism
$f\colon B \to B'$ such that the given maps $A \to B'$ and
$Y_2 \to B(k)$ can be obtained by composing those of $B$ with $f$.
%% NOTE (p52): "and $Y_2 \to B(k)$" above is as printed (no prime on $B$;
%% the parallel clause earlier in the same sentence prints $Y_2 \to B'(k)$).
(In other words, $B$ is a \emph{push-out} over $Y_1$ of $A$ and the
``constant'' group scheme defined by $Y_2$.)

The uniqueness of $B$ is assured by its universality.  Let us prove its
existence.  For each $y \in Y_3$ let $\bar{y}$ be a representative of $y$
in $Y_2$.  If $y, y' \in Y_3$, we have
\[
\bar{y} + \bar{y'} = \overline{y + y'} + c(y, y')
\]
with $c(y, y') \in Y_1$; the cochain $c$ is a $2$-cocycle defining the
extension \eqref{eq:II.1.3.star}.  Let $B$ be the disjoint union of copies
$A_y$ of $A$, indexed by $y \in Y_3$.  Define a group law on $B$ via the
mappings
\[
\pi_{y,y'} : A_y \times A_{y'} \longrightarrow A_{y+y'}
\qquad (y, y' \in Y_3),
\]
given by addition in $A$ followed by translation by $\epsilon(c(y,y'))$.
One then checks easily that $B$ has the required universal property, the
maps $A \to B$ and $Y_2 \to B(k)$ being defined as follows:

$A \to B$ is the natural map $A \to A_o$ followed by translation by
$-c(0, 0)$,
%% NOTE (p53): the subscript of $A$ above is printed as a small round
%% letter "o", clearly distinct from the narrow oval digit "0" of the
%% same page (cf. "$-c(0, 0)$", "0 -> A(k')"); transcribed literally
%% as $A_o$ (same convention as $K_o$ in chap. I, 1.2).

$Y_2 \to B(k)$ maps an element $\bar{y} + z$, $y \in Y_3$, $z \in Y_1$
onto the image of $z$ in $A_y$.

Note that for any extension \emph{field} $k'$ of $k$ we have an
\emph{exact sequence}
\[
0 \longrightarrow A(k') \longrightarrow B(k') \longrightarrow Y_3
\longrightarrow 0,
\]
and a \emph{commutative diagram}
\[
\begin{array}{ccccccccc}
0 & \to & Y_1 & \to & Y_2 & \to & Y_3 & \to & 0 \\
  &     & \downarrow & & \downarrow & & \downarrow & & \\
0 & \to & A(k') & \to & B(k') & \to & Y_3 & \to & 0.
\end{array}
\]
The algebraic group $B$ is thus an \emph{extension} of the ``constant''
algebraic group $Y_3$ by $A$.

\begin{remarks}
\begin{nlist}
\item Let $k'$ be an extension of $k$ and $A' = A \times_k k'$.  We may
apply the above construction to the $k'$-algebraic group $A'$, with
respect to the exact sequence \eqref{eq:II.1.3.star} and to the map
$Y_1 \to A(k) \to A'(k')$.  The group $B'$ thus obtained is canonically
isomorphic to $B \times_k k'$; this follows, for instance, from the
explicit construction of $B$ and $B'$.

\item We will only use the above construction when $\mathrm{char}(k) = 0$
and $A$ is a torus.  The enlarged group $B$ is then a ``group of
multiplicative type''; this means that, after a suitable finite extension
of the ground field, $B$ becomes isomorphic to the product of a torus and
a finite abelian group.  Such a group is uniquely determined by its
character group $X(B) = \Hom_{\bar{k}}(B_{/\bar{k}},
\mathbf{G}_{m/\bar{k}})$, which is a Galois-module of finite type over
$\Z$.  Here $X(B)$ can be described as the set of pairs $(\phi, \chi)$,
where $\phi\colon Y_2 \to \bar{k}^*$ is a homomorphism and
$\chi \in X(A)$ is such that $\phi(y_1) = \chi(y_1)$ for all
$y_1 \in Y_1$.  Note that this gives an alternate definition of $B$.
\end{nlist}
\end{remarks}

\runhead{Exercise}
\begin{alist}
\item[a)] Let $k'$ be a commutative $k$-algebra, with $k' \neq 0$, and
$\Spec(k')$ connected (i.e.~$k'$ contains exactly two idempotents: $0$
and $1$).  Show the existence of an exact sequence:
\[
0 \longrightarrow A(k') \longrightarrow B(k') \longrightarrow Y_3
\longrightarrow 0
\]
\item[b)] What happens when $\Spec(k')$ is not connected?
\end{alist}
